% vim: ai sw=2 spell spelllang=cs

\begin{frame}
  \frametitle{Datové typy}

  \heading{LDD3 kap. 11}

  \bigskip

  \begin{itemize}
  \item Základní typy
%  \item Přenositelnost
  \item Endianita (pořadí bytů)
  \item Datové struktury (seznam, FIFO)
  \end{itemize}

\end{frame}

\mysection{Základní typy}

\begin{frame}
  \frametitle{Základní typy}

  \begin{itemize}
  \item \code{char}, \code{short}, \code{int}, \code{long},
    \code{long long} (+ \code{unsigned})
  \addtocounter{beamerpauses}{1}
  \item<+-> \code{u8}, \code{u16}, \code{u32}, \code{u64}\uncover<+->{,
    \code{__u8}, \code{__u16}, \code{__u32}, \code{__u64}}
  \item<.(-1)-> \code{s8}, \code{s16}, \code{s32}, \code{s64}\uncover<.->{,
    \code{__s8}, \code{__s16}, \code{__s32}, \code{__s64}}
  \end{itemize}

  \addtocounter{beamerpauses}{-1}
  \pause

  \begin{block}{\code{sizeof} základních typů pro různé architektury}
  \begin{center}
  \footnotesize
  \begin{tabular}{|l|c|c|c|c|c|c|} \hline
Architektura & \code{char} & \code{short} & \code{int} & \code{long} &
  \code{void *} & \code{long long} \\ \hline\hline
alpha   & 1    & 2     & 4   & 8    & 8   & 8 \\ \hline
arm     & 1    & 2     & 4   & 4    & 4   & 8 \\ \hline
ia64    & 1    & 2     & 4   & 8    & 8   & 8 \\ \hline
m68k    & 1    & 2     & 4   & 4    & 4   & 8 \\ \hline
mips    & 1    & 2     & 4   & 4    & 4   & 8 \\ \hline
ppc     & 1    & 2     & 4   & 4    & 4   & 8 \\ \hline
sparc   & 1    & 2     & 4   & 4    & 4   & 8 \\ \hline
sparc64 & 1    & 2     & 4   & 8    & 8   & 8 \\ \hline
x86\_32 & 1    & 2     & 4   & 4    & 4   & 8 \\ \hline
x86\_64 & 1    & 2     & 4   & 8    & 8   & 8 \\ \hline
  \end{tabular}
  \end{center}
  \end{block}

\end{frame}

\begin{frame}
  \frametitle{Základní typy}

  \begin{block}{\code{sizeof} dalších typů (ukazatelů)}
  \begin{center}
  \footnotesize
  \begin{tabular}{|l|c|c|c|} \hline
Architektura & \code{long} & \code{void *} & Fyzická \\ \hline\hline
x86\_32     & 4    & 4    & 4    \\ \hline
x86\_32+PAE & 4    & 4    & 4.5+  \\ \hline
x86\_64     & 8    & 8    & 8    \\ \hline
  \end{tabular}
  \end{center}
  \end{block}

  \pause
  \medskip

  \heading{Pro práci s fyzickou adresou \texttt{long} nebo \texttt{void *}
    nestačí}
  \begin{itemize}
    \item Ale \code{u64} nebo \code{long long} může být zbytečné (čisté
      32-bity)

    \pause

    \item Proto speciální typy (\emph{POZOR})
    \begin{itemize}
      \item \code{phys_addr_t} (RAM, za MMU)
      \item Anebo \code{long}, ale pak v něm uchovávat jen \code{pfn}
	(\code{phys/PAGE_SIZE})

      \pause

      \item \code{dma_addr_t} (RAM, před IOMMU)
      \item \code{resource_size_t} (PCI prostor apod.)
    \end{itemize}
  \end{itemize}

  \pause
  \medskip

  \heading{Úkol:} opravte typy v pb173/06 (jen \emph{první} TODO)

\end{frame}

\begin{frame}[fragile]
  \frametitle{Pořadí bytů}

  \begin{block}{Číslo \code{0x12345678} reprezentované v paměti}
  \begin{center}
   \begin{tabular}{l|l|llll}
\multirow{2}{*}{\textbf{Pořadí}} & \textbf{Příklad} &
  \multicolumn{4}{c}{\textbf{Offset v paměti}} \\
              & \textbf{architektury}       &    0 &    1 &    2 &    3 \\\hline
Big-endian    & \texttt{arm}, \texttt{s390} & 0x12 & 0x34 & 0x56 & 0x78 \\
Little-endian & \texttt{arm}, \xes          & 0x78 & 0x56 & 0x34 & 0x12 \\
Mixed-endian  & \texttt{PDP-11}             & 0x34 & 0x12 & 0x78 & 0x56 \\
  \end{tabular}
  \end{center}
  \end{block}

  \pause
  \medskip

  \heading{Pro typy o velikosti větší než 1 je nutné znát pořadí}
  \begin{itemize}
    \item Interpretace dat z/do zařízení

    \pause

    \item Specifikace
    \begin{itemize}
      \item CPU (příklady nahoře)
      \item PCI: LE
      \item Síťový provoz: BE
      \item \dots
    \end{itemize}
  \end{itemize}

\end{frame}

\begin{frame}[fragile]{Funkce pro pořadí bytů}

  \begin{itemize}
    \item \header{asm/byteorder.h}
    \item \code{__leXX}, \code{__beXX}
    \item $\text{\code{XX}}\in\left\{\text{\code{16}},\text{\code{32}},\text{\code{64}}\right\}$

    \pause

    \item \code{cpu_to_leXX}, \code{cpu_to_beXX}
    \item \code{leXX_to_cpu}, \code{beXX_to_cpu}

    \pause

    \item \code{ntohs}, \code{ntohl}, \code{htons}, \code{htonl}
  \end{itemize}

  \pause
  \bigskip

  \heading{Úkol:} doplňte tělo \code{decode_and_print} v pb173/06 (druhé TODO)

  \onslide<2->

  \begin{block}{Příklad}
    \begin{lstlisting}
__le32 my_le32 = cpu_to_le32(0x1234);
u32 my_u32 = le32_to_cpu(my_le32);
    \end{lstlisting}
  \end{block}

\end{frame}

\iffimuni
\subsection{Chyby jako ukazatele}
\fi

\begin{frame}
  \frametitle{Chyby jako ukazatele}

  \begin{itemize}
    \item \code{void *funkce()}
    \item Jak vrátit chytřejší chybu než NULL?

    \pause

    \begin{itemize}
      \item V uživatelském prostoru máme \code{errno} a ekvivalenty
    \end{itemize}

    \pause

    \item Díra na konci adresového prostoru
    \begin{itemize}
      \item Tj.\ nevalidní ukazatel (dereference by způsobila pád)
      \item Např.\ na \xesh: \code{0xfffffffffff00000-0xffffffffffffffff}
      \item Lze použít k zakódování chyby
    \end{itemize}

    \pause

    \item \code{void *ptr = ERR_PTR(-Eerror)}
    \item \code{IS_ERR(ptr) == true} $\Rightarrow$ \code{int err = PTR_ERR(ptr)}

    \pause

    \item \code{IS_ERR(NULL)} je \code{false}!
  \end{itemize}

  \pause
  \bigskip

  \heading{Úkol:} přepište \texttt{error\_handling} v pb173/06 (třetí TODO)

\end{frame}

\mysection{Datové struktury}

\iffimuni
\subsection{Seznam}
\fi

\begin{frame}[fragile]
  \frametitle{Seznam I.}

  \begin{itemize}
    \item \header{linux/list.h}
    \item \code{struct list_head} (obsahuje \code{prev}, \code{next})
    \begin{itemize}
      \item \emph{Jak samotný seznam (počátek), tak jeho prvky}
    \end{itemize}

    \pause

    \item \code{LIST_HEAD(my_list)}$_S$, \code{INIT_LIST_HEAD(\&my_list)}$_D$

    \pause

    \item \code{list_add(co, kam)}, \code{list_add_tail(co, kam)},
      \code{list_del(co)}
    \item \code{list_move(co, kam)}, \code{list_move_tail(co, kam)}
  \end{itemize}

  \pause

  \begin{block}{Příklad}
  \begin{tabular}{lp{.03\textwidth}l}
  \begin{lstlisting}[language=C]
struct my_struct {
  int a;
  struct list_head list;
};
LIST_HEAD(my_list);
  \end{lstlisting} & &
  \begin{lstlisting}
struct my_struct *s;
s = kmalloc(sizeof(*s), GFP_KERNEL);
if (!s) { ... }
s->a = 10;       /* or list_add_tail(&s->list, &my_list); */
list_add(&s->list, &my_list);
  \end{lstlisting}
  \end{tabular}
  \end{block}

  \pause
  \medskip

  \heading{Úkol:} naalokujte 20 stránek a v seznamu si je pamatujte.

\end{frame}

\begin{frame}[fragile]{Seznam II.}

  \begin{itemize}
    \item \header{linux/list.h}
    \item \code{list_empty(list)}

    \pause

    \item Jedna položka: \code{list_entry}, \code{list_first_entry}

    \pause

    \item Iterace:
      \code{list_for_each_entry(...) \{ ... \}} (jako \code{for}), \\
      \code{list_for_each_entry_reverse(...) \{ ... \}}

    \pause

    \item \code{*_safe} varianty -- pokud měním aktuální člen seznamu
  \end{itemize}

  \onslide<2->

  \begin{block}{Příklad}
    \begin{lstlisting}[aboveskip=2pt]
struct my_struct *s, *s1;
s = list_first_entry(&my_list, struct my_struct, list);
s = list_entry(s->list.next, struct my_struct, list);
    \end{lstlisting}
    \onslide<3->
    \begin{lstlisting}[belowskip=0pt]
list_for_each_entry(s, &my_list, list) { s->a; }
    \end{lstlisting}
    \onslide<4->
    \begin{lstlisting}[aboveskip=0pt,belowskip=2pt]
list_for_each_entry_safe(s, s1, &my_list, list) { list_del(&s->list); }
    \end{lstlisting}
  \end{block}

  \pause
  \medskip

  \heading{Úkol:} předchozích 20 stránek uvolněte.

\end{frame}

\iffimuni
\subsection{FIFO}
\fi

\begin{frame}{FIFO}

  \begin{itemize}
    \item \header{linux/kfifo.h}, \doc{samples/kfifo/*}
    \item \code{struct kfifo}
    \begin{itemize}
      \item Statické: \code{DECLARE_KFIFO} + \code{INIT_KFIFO}
      \item Dynamické: \code{kfifo_alloc(\&fifo, size, GFP_*)},
	\code{kfifo_free(&fifo)}
    \end{itemize}

    \pause

    \item Zápis: \code{kfifo_in(\&fifo, buf, count)}
    \item Čtení: \code{kfifo_out(\&fifo, buf, count)}

    \pause

    \item Obsazeno: \code{kfifo_len(\&fifo)}
    \item Volno: \code{kfifo_avail(\&fifo)}
  \end{itemize}

\end{frame}

\begin{frame}{Úkol}

  \heading{Práce s FIFO} (\header{linux/kfifo.h})
  \begin{enumerate}
    \item Globální FIFO 8 intů (\code{DECLARE_KFIFO})
    \item \code{INIT_KFIFO}
    \item Zapisovat čísla (\code{kfifo_in}) dokud je místo
    \item Vypsat všechna čísla (\code{kfifo_out} + \code{kfifo_len})
  \end{enumerate}

  \bigskip

  Pozn.: ve starších jádrech ($<$ 2.6.33) je jiné rozhraní:
  \begin{itemize}
    \item jen dynamicky: \code{my_fifo = kfifo_alloc}
    \item \code{__kfifo_put} je \code{kfifo_in}
    \item \code{__kfifo_get} je \code{kfifo_out}
    \item \code{__kfifo_len} je \code{kfifo_len}
    \item \code{avail} není
  \end{itemize}

\end{frame}

